\subsubsection{Patching at an absolutely irreducible ordinary arc}
\label{par:ordinary-patching}

The finite-level transfer controls the trace relations before passing to a limit. Patching uses that control to deduce ordinary support along an arc when its generic representation satisfies the stated irreducibility and local separation conditions.

\begin{proposition}[Augmentation support below a marked arc]\label{prop:ordinary-one-point-patching}
In one rank-four determinant link, take a characteristic-\(\ell\)
ordinary arc through its ordinary support whose representation is
absolutely irreducible on the cyclotomic kernel of \(G_L\).
Assume the ordered characters separate on the prescribed inertia
groups, regular elements are dense in the cyclotomic kernel, and
the distinguished auxiliary place satisfies the local cyclotomic
intertwining condition of Proposition~\ref{prop:tower-below-arc}.
Assume \(\bar\epsilon(G_L)\ne\{1\}\), where \(\bar\epsilon\)
is the mod-\(\ell\) cyclotomic character.
With the finite-level modules and auxiliary comparisons of
Propositions~\ref{prop:ordinary-level-modules}
and \ref{prop:ordinary-integral-transfer}, if \(y\) denotes the arc and \(R_0\) the base ring of the Hecke ideal, then every generization satisfies 
\[
 P\subseteq\ker y\quad\Longrightarrow\quad J_0\subseteq P\cap R_0.
\]

\end{proposition}

\begin{proof}
\proofstep{One determinant link and its support condition.}

Fix one of the two determinant links and write \(R_0\) for its
universal ordinary determinant ring, \(J_0\) for its Hecke ideal,
and \(p\) for the evaluation kernel of the characteristic-\(\ell\)
arc. Write \(n=4\), \(d=[L^+:\Q]\), and
\[
 D=nd.
\]
We use the
one-point case in which \(p\) contains \(J_0\) and the arc
representation is absolutely irreducible on the cyclotomic kernel
of \(G_L\). Use the ordered local characters and auxiliary condition
of Proposition~\ref{prop:tower-below-arc}
and the residual tests of Lemma~\ref{hyp:ordinary-arc-residual-tests}. We describe
the passage from finite-level data to vanishing of \(J_0\) at the
generizations of \(p\).

Put
\[
 A=K[[t]]
\]
for the arc valuation ring after finite
extensions and \(E=K((t))\). The finite field \(K\) can be taken
to contain the coefficient residue field. Choose Taylor--Wiles
elements \(\sigma_j\), \(1\le j\le q\), in the full cyclotomic
kernel of \(G_L\) at which the characteristic polynomials are
separable. Here \(q\) counts the Taylor--Wiles places and is
fixed throughout the patching limit. The elements must detect
dual adjoint classes, as follows. On the cyclotomic kernel the
adjoint action twisted by the cyclotomic character becomes the
ordinary adjoint action. Write
\[
 V(g)=z(g)\rho(g)
\]
for a dual
cocycle \(z\) on that kernel. Its cocycle identity is
\[
 V(gh)=V(g)\rho(h)+\rho(g)V(h).
\]
Suppose \(\operatorname{tr} V(g)=0\) for every \(g\) in the
kernel.

If \(\sum_i a_i\rho(g_i)=0\), the lift identity
\(V(gh)=V(g)\rho(h)+\rho(g)V(h)\) gives
\[
 \operatorname{tr}\left(\Bigl(\sum_i a_i V(g_i)\Bigr)\rho(h)\right)
  = \sum_i a_i\operatorname{tr}V(g_i h)
        -\operatorname{tr}\left(\Bigl(\sum_i a_i\rho(g_i)\Bigr)V(h)\right)
  =0.
\]
Every \(g_i h\) here lies in the same cyclotomic kernel,
where the traces are assumed zero.
Absolute irreducibility gives the full span \(M_n(E)\).
Its trace pairing is perfect, so \(\sum_i a_i V(g_i)=0\).
Thus the proposed map on the spanning matrices is well-defined
and obeys the Leibniz rule.

Write \(D_\rho\) for this derivation and
\(e_{ij}\) for the matrix units. The matrix
\(Y=\sum_i D_\rho(e_{i1})e_{1i}\) implements the derivation.
Indeed the Leibniz rule applied to \(e_{ab}e_{i1}=\delta_{bi}e_{a1}\) gives
\[
 Ye_{ab}=D_\rho(e_{a1})e_{1b},\qquad
 e_{ab}Y=D_\rho(e_{a1})e_{1b}-D_\rho(e_{ab}).
\]
Hence \(V(g)=Y\rho(g)-\rho(g)Y\) on the cyclotomic kernel \(H\).
Subtracting the associated coboundary makes \(z|_H=0\).

For \(h\in H\) and \(g\in G_L\), apply the cocycle identity to
\(hg=g(g^{-1}hg)\). Normality and \(z|_H=0\) give
\[
 \rho(h)z(g)\rho(h)^{-1}=z(hg)=z(g).
\]
The span of \(\rho(H)\) is the full matrix algebra.
Its commutant is \(E\cdot I_n\), so \(z\) descends to a scalar
cocycle on the abelian quotient \(G_L/H\), with the residual
cyclotomic action \(\epsilon=\bar\epsilon\).
For the fields used here a place with nontrivial residual cyclotomic
value is provided by the marked-field constructions in
Propositions~\ref{prop:count-ii-auxiliary-primes}
and~\ref{prop:count-three-marked-elements}. Choose \(a\) in
the quotient with \(\epsilon(a)\ne1\) and put
\(u=z(a)/(\epsilon(a)-1)\). Since \(ab=ba\),
\[
 z(a)+\epsilon(a)z(b)=z(b)+\epsilon(b)z(a)
 \quad\Longrightarrow\quad z(b)=(\epsilon(b)-1)u .
\]
This exhibits the coboundary.

Passing from the real group to its index-two CM subgroup
loses no classes, since \(2\) is invertible. The regular
elements are dense in the kernel, so continuity suffices to
test traces on that locus. The space of Laurent series classes
is finite-dimensional by finite generation after passage to
the pro-\(\ell\) quotient. Choose a finite set of its trace
tests. After an extension of \(A\), each marked characteristic polynomial takes the form
\[
 P_{\sigma_j}(X)=\prod_{a=1}^n (X-\alpha_{ja}),\qquad
 \alpha_{ja}\in A,\qquad \prod_{a<b}(\alpha_{ja}-\alpha_{jb})\ne0\text{ in }E .
\]
The differences need only be nonzero in the fraction field here.

For every depth \(N\) choose distinct split places away from
the initial level with Frobenius approaching \(\sigma_j\) and
norm
\[
 Q_{j,N}\equiv1\pmod{\ell^N}.
\]

For precision \(m\), scale the fixed finite list of integral
cocycles once and record their reductions alongside the residual
matrix action modulo \(t^m\). For a cocycle \(z\), conjugation on its coefficient module gives the homomorphism
\[
 h\longmapsto (z(h),r(h))\bmod t^m,\qquad
 (u,A)(v,B)=(u+\operatorname{Ad}(A)v,AB).
\]

The coefficient and matrix rings modulo \(t^m\) are finite. Intersect their open kernels
with the kernel of the cyclotomic character modulo \(\ell^N\);
take the normal core in the group over the totally real field,
also recording its CM quotient. All these quotients are finite.
Use \(m\ge N\) tending to infinity in the resulting sequence.
Chebotarev gives infinitely many primes in each specified
conjugacy class, so we may avoid all previous choices and the
fixed bad primes.
In the resulting finite quotient \(G_{m,N}\) the target class is
\[
          [\sigma_j]_{G_{m,N}},\qquad
          \epsilon_\ell(\sigma_j)\equiv1\bmod\ell^N .
\]
The split condition comes from its membership in the CM
subgroup. Use the inverse conjugacy class for geometric Frobenius.
Choose the tame generators in the pro-\(\ell\) inertia factors. The open subgroup on which all
residual characters and the mod-\(\ell\) cyclotomic character
are trivial has pro-\(\ell\) image for the arc.

In fact its
residual matrices are unipotent, so the finite residual
image is an \(\ell\)-group by Cauchy's theorem. At depth \(m\) the kernels are filtered by the groups
\(1+t^aM_n(K[t]/(t^m))\), with successive quotients embedded in additive \(M_n(K)\) by
\[
 1+t^aX\mapsto X\bmod t.
\]
Indeed
\[
 (1+t^aX)(1+t^aY)\equiv1+t^a(X+Y)\pmod{t^{a+1}}
 \qquad (a\ge1).
\] Thus each is an \(\ell\)-group. This justifies the use of the
maximal pro-\(\ell\) quotient on this subgroup. The finite
levels are compared on a common topologically finitely
generated profinite group, with stable quotients in every
fixed finite approximation along an ultrafilter
\(\mathcal U\) containing the cofinite filter.
Frobenius and tame inertia converge in these quotients.

To obtain the common group,
bound the generator ranks of the open pro-\(\ell\)
subgroups uniformly. A homomorphism to \(\mathbf F_\ell\)
which kills all lifts of the new tame generators descends
to the base. If \(d_0\) is the base generator rank, \(q\) the number
of new places, and \(I_0\) the fixed open subgroup index,
the evaluations on at most \(qI_0\) such lifts give
\[
 \dim_{\mathbf F_\ell}\operatorname{Hom}_{\rm cont}
      (U_N,\mathbf F_\ell)\le d_0+qI_0 .
\]
The Frattini quotient identifies the left side with the
generator rank of the pro-\(\ell\) group \(U_N\).
Use a free profinite group on this many generators
together with lifts of the fixed finite quotient.

In each finite quotient of the common group there are
finitely many possibilities for the image of the level
kernel. The ultrafilter selects compatible stabilized
images. Their inverse limit is the kernel of the limiting
quotient. The decomposition maps are carried along in
this construction. The ring \(R\) below uses the
determinant on the common group, and \(I\) imposes
its descent to the limiting quotient.

Use a coefficient DVR \(O\) with residue field \(K\) after
finite unramified coefficient extension, and adjoin a variable
\(s\) on the base, evaluated at \(t\). These extensions are
flat. The completed ordinary weight algebra on the base is
\[
 \Lambda_0=O[[s,w_1,\ldots,w_D]].
\]
Its specialization to \(K[[t]]\) is onto because it sends
\(s\) to \(t\). The completion of \(\Lambda_0\) at the
kernel of this specialization is regular local of dimension
\(1+D\). Indeed the quotient has dimension one, and
\(\Lambda_0\) is regular local of dimension \(2+D\).

At the finite levels retain one tame character variable
\(y_{ji}\) for each Frobenius eigenline at each of the \(q\)
auxiliary places. Put \(\mathfrak y=(y_{ji})\). The required
finite-level ordinary module \(M_N\) is finite free over
\[
 \Lambda_0[[y_{ji}]]\Big/
 \bigl((1+y_{ji})^{\ell^{a_{j,N}}}-1\bigr)_{j,i},
 \qquad a_{j,N}\ge N .
\]

Adjoin the ordered roots of the characteristic polynomials at
the \(\sigma_j\) over the extended base ring, using one monic
polynomial and then its successive monic quotients. If \(P_0\) is the polynomial of degree \(n\) at one place, the successive extensions are
\[
 C_i=C_{i-1}[z_i]/P_{i-1}(z_i),\qquad
 P_i(X)=P_{i-1}(X)/(X-z_i),\qquad
 \operatorname{rank}_{C_{i-1}} C_i=n+1-i .
\]

Thus ordering the roots is finite free, including when the residual
polynomial has repeated roots. If the residual labeling singles out a
maximal ideal, take the corresponding direct factor \(R_0'\):
over a complete local base the idempotents of the finite special
fibre lift uniquely by Hensel's lemma. This summand is finite
projective, hence free over the local base. If \(C\) is the
finite splitting algebra and \(\bar e\) the residual projector, this
construction gives
\[
             C=eC\oplus(1-e)C,\qquad e^2=e,\quad e\bmod\mathfrak m=\bar e .
\]
Evaluation at the
chosen roots selects this summand and lifts the arc. Differences
which vanish at its closed point can still be nonzero at the
generic point of the arc. The base module \(X\) is the ordinary
Hecke module after this base change. It is finite free
over \(\Lambda_0\); its support is the inverse image of
\(V(J_0)\).

At a finite level use the roots \(U_{ji}\) of
the corresponding Frobenius. The maps between the
zero-\(\mathfrak y\) fibre of \(M_N\) and the base module
with those ordered roots have both compositions equal to
multiplication by
\[
 f_N=\prod_j\prod_{i}\prod_{i'\ne i}
 (U_{ji}-U_{ji'})(U_{ji}-Q_{j,N}U_{ji'}).
\]
These are the Hecke-module freeness and level comparison
inputs. The second factor is needed in the comparison
at finite depth. At the limiting arc both factors become
the same root difference, which is nonzero in \(E\).
The integral construction at these levels is described
in Paragraph~\ref{par:ordinary-level}. It gives the
trace at the tame cover, the two comparisons and the
uniform exponent for the local relations.

Here is the algebraic form of the remaining finite-level
compatibilities. One has one complete Noetherian local
ring \(R\) acting on the modules, with coefficient map
\[
 \Lambda_0[[y_{ji}]]\to R.
\]
The coefficients of the determinant of the limiting group supply the action of \(R\).
There is a surjection \(R\to R_0'\) killing \(\mathfrak y\),
and an ideal \(I\) in its kernel which imposes the limiting
ordinary, polarized, and prescribed inertia relations.
The Frobenius and tame trace relations use the following
elements in the group algebra or its Cayley--Hamilton quotient:
\[
       \prod_i(\sigma_j-U_{ji}),\qquad
       \Big(\prod_{i'\ne i}(\sigma_j-U_{ji'})\Big)
                       (\tau_j-1-y_{ji}).
\]
At the limiting point the marked Frobenius has simple spectrum. Its commutation with the tame element implies that each Frobenius line \(L_{ji}\) is tame-stable. The second relation then reads
\[
 \rho(\tau_j)|_{L_{ji}}=(1+y_{ji})\,\mathrm{id}_{L_{ji}}
\]
after the distinct-root differences have been inverted. Pairing these
relations with the traces of all group elements gives
coefficient relations.
\proofstep{Comparison element and patching depth.}

For the proof below the comparisons give an element
\(f\in R\) specializing to the product of the squared
root differences and the following two properties.
\begin{itemize}
\item For every fixed \(b\), for \(\mathcal U\)-almost
all \(N\) there are \(R\)-linear maps in both directions
between \(X/\mathfrak m_{\Lambda_0}^bX\) and
\[
 M_N/(\mathfrak y,\mathfrak m_{\Lambda_0}^b)M_N,
\]

with both composites multiplication by \(f\).
\item For a generator \(r\) of \(I\), some fixed power
\((fr)^e\) agrees modulo each fixed power of \(\mathfrak m_R\)
with an element annihilating \(M_N\) for
\(\mathcal U\)-almost all \(N\). The exponent is chosen
before the accuracy and depth.
\end{itemize}
The second property permits congruences of the action at
different finite levels. Here the generators of \(I\) are the trace relations for the group kernel,
polarization, ordinarity, the prescribed inertia, and
the two Frobenius and tame relations just described.

We prove the patching comparison for these data.
For \(k,b\ge1\) let \(V_{k,b}\) be the ultraproduct
along \(\mathcal U\) of
\[
 M_N/(\mathfrak y^k,\mathfrak m_{\Lambda_0}^b)M_N.
\]

For fixed \(k,b\) the group-ring relations are in
\((\mathfrak y^k,\mathfrak m_{\Lambda_0}^b)\) for all
large \(N\). For \(1\le c<k\le\ell^a\), the identity
\[
 c\binom{\ell^a}{c}=\ell^a\binom{\ell^a-1}{c-1}
 \quad\Longrightarrow\quad
 v_\ell\binom{\ell^a}{c}\ge a-v_\ell(c)
\]
is uniform in these finitely many \(c\) as \(a\) tends to infinity.
\proofstep{Compatible bases and the augmentation fibre.}

Choose a free basis of \(M_N\) over its coefficient group
ring at each \(N\). For fixed \(k,b\) the binomial valuation
bound kills all relations of positive \(y\)-degree less than
\(k\) modulo \(\mathfrak m_{\Lambda_0}^b\) for large \(N\).
Writing \(d_N\) for the size of the chosen basis, the
basis therefore gives the additive decomposition
\[
 M_N/(\mathfrak y^k,\mathfrak m_{\Lambda_0}^b)M_N
   \simeq\bigoplus_{|\alpha|<k}
       (\Lambda_0/\mathfrak m_{\Lambda_0}^b)^{d_N}y^\alpha .
\]
Reduction in \(b\) acts on the coefficients and reduction
in \(k\) discards the extra monomials. Thus both commute with
the identifications from the fixed bases. For fixed \(k\), take the inverse limit
over \(b\), giving \(V_k\). Multiplication by the monomials
identifies the successive graded pieces with copies of
\(V_1\). These identifications are \(R\)-linear, because
multiplication by each \(y_{ji}\) commutes with the
\(R\)-action.

At \(k=1\), the comparison maps give the ultraproduct and inverse
limit maps
\(X \xrightarrow{u} V_1\xrightarrow{v} X\).
To arrange compatibility in \(b\), choose maps at a depth
of truncation tending to infinity along the ultrafilter
and use their reductions. Indeed, if \(S_j\in\mathcal U\) is the good set for truncation \(j\), then
\[
 T_b=\bigcap_{j=1}^b S_j\in\mathcal U,\qquad T_{b+1}\subseteq T_b .
\]
On \(T_b\) maps at depth \(b\) can be reduced compatibly to every smaller truncation. At level \(N\) select the largest
available truncation up to \(N\), which gives the
required cofinal choices. The
ultraproduct of the constant finite module
\(X/\mathfrak m_{\Lambda_0}^bX\) is that same module.
Their composites satisfy
\[
      vu=f\,\mathrm{id}_X,\qquad uv=f\,\mathrm{id}_{V_1}.
\]
Inverting \(f\) gives \(V_1[1/f]\simeq X[1/f]\).

Now form
\[
 P=\varprojlim_k V_k[1/f]
\]
over the \(\mathfrak y\)-adic completion of \(R[1/f]\).
Flatness of localization preserves the transition kernels. The
monomial expansions identify
\[
   \operatorname{gr}_{\mathfrak y}P\simeq X[1/f][y_{ji}],
           \qquad P/\mathfrak yP\simeq X[1/f].
\]
The degree-\(a\) identification is multiplication by all monomials
of total degree \(a\).
For this assertion, a compatible sequence with zero
constant term is a sum of the \(y_{ji}\) times compatible
sequences. Solve this degree by degree with the
monomial identifications. The same argument applies in
each degree. Finite generators of \(X[1/f]\) lift to
generators of \(P\) by lifting the coefficient at each
successive degree and summing in the complete ring.

The relations in \(I\) act nilpotently on \(P\). To see
why congruences in the maximal ideal of \(R\) suffice,
fix \(b,k\). Some \(\mathfrak m_R^c\) kills
\(X/\mathfrak m_{\Lambda_0}^b X\), a finite module over
the local ring \(R\). The comparison at this truncation
therefore gives \(f\mathfrak m_R^cV_{1,b}=0\).
Iterating on successive \(\mathfrak y\)-quotients gives
\[
 f^k\mathfrak m_R^{ck}V_{k,b}=0.
\]
Use the annihilator approximation modulo
\(\mathfrak m_R^{ck}\). For every \(k,b\), the bound gives \(f^k(fr)^eV_{k,b}=0\).
After taking the inverse limit in \(b\) and inverting \(f\),
\[
             r^e V_k[1/f]=0 .
\]
Since \(e\) is independent of \(k\), for each compatible sequence \(v=(v_k)_k\in P\) we have
\[
 r^e v=(r^ev_k)_k=(0)_k.
\]
This gives \(r^e P=0\) on the inverse limit.

Let \(p_R\) be the inverse image in \(R\) of the lifted
arc prime in \(R_0'\), and put
\[
 C=\widehat{R_{p_R}}.
\]
Its residue field is \(E\). At each quotient by \(p_R^m\), the
inclusion \(\mathfrak y^m\subseteq p_R^m\) makes the preliminary
\(\mathfrak y\)-completion redundant, and \(f\notin p_R\).
Taking the localized inverse limit gives the same ring \(C\).
This base change on the intermediate Noetherian ring is
flat. Write \(P_C\) for the completed local base change
of \(P\), and \(\Lambda_p\) for the regular completed
localization of \(\Lambda_0\) above. The augmentation
\(P_C/\mathfrak y P_C\) is finite free over
\(\Lambda_p\).

To see this, the action on \(X\) factors
through a finite \(\Lambda_0\)-algebra inside
\(\operatorname{End}_{\Lambda_0}(X)\). After base change to the complete local ring \(\Lambda_p\), write the finite algebra as \(\prod_{\mathfrak q}T_{\mathfrak q}\) and its idempotents as \(e_{\mathfrak q}\). Then
\[
 X\otimes\Lambda_p=\bigoplus_{\mathfrak q} e_{\mathfrak q}(X\otimes\Lambda_p).
\]
The summand for \(p_R\) is projective over the local ring \(\Lambda_p\), hence finite free. This factor is
nonzero by \(J_0\subset p\). Flat base change preserves
the monomial identifications.

If \(P_C/\mathfrak yP_C\) has rank \(h\) over \(\Lambda_p\),
lift a basis. Successive monomial coefficients give the surjection
\(\Lambda_p[[y_{ji}]]^h\to P_C\). Indeed, put \(S=\Lambda_p[[y_{ji}]]\) and let \(e_i\) be the chosen lifts. A nonzero relation would have a least total degree \(a\). Thus its initial part would give
\[
 \sum_{i,\,|\alpha|=a} c_{i,\alpha}\,y^\alpha\bar e_i=0
 \quad\text{in }\mathfrak y^aP_C/\mathfrak y^{a+1}P_C .
\]
The monomial identification makes these vectors a \(\Lambda_p\)-basis, so all the initial coefficients vanish. Hence \(S^h\simeq P_C\). Completeness of \(C\) and finiteness
of \(P_C\) give completeness for the \(\mathfrak y\)-adic
lifting.

The support algebra \(C/\operatorname{Ann}_C P_C\)
embeds into the endomorphism algebra of this finite free
module over \(\Lambda_p[[y_{ji}]]\). It is finite and
torsion-free over that regular local ring. Every minimal
prime of the support algebra contracts to zero, since
each nonzero element of the regular ring is a
nonzerodivisor and minimal primes are associated.
Dimension for the finite integral extension now shows
that every minimal support prime \(Z\) has
\[
 \dim(C/Z)=1+D+nq .
\]
Each \(Z\) contains \(IC\), by the annihilator
calculation above. For \(Q\supseteq\mathfrak y\), Nakayama at \(Q\) gives
\[
    (P_C)_Q\ne0\quad\Longleftrightarrow\quad
                      (P_C/\mathfrak yP_C)_Q\ne0 .
\] In terms of \(R_0'\) it is
therefore equivalent to containment of \(J_0\) in
the contraction.

We explain the upper bound on the relative cotangent
dimension of \(C/IC\) needed to use the support.
The coefficient DVR at this point can be taken to be
the completed localization of \(O[[s]]\) at the
coefficient uniformizer. Its residue field is \(E\),
because \(s\) specializes to \(t\).
If \(T_{\rm rel}\) is the tangent space relative to this DVR, the first-order estimate to prove is
\[
 \dim_E T_{\rm rel}\le D+nq.
\]
Here the rank-four point has one
absolutely irreducible matrix block. The localized
Cayley--Hamilton algebra is the full matrix algebra
\cite[Theorem 2.22(i)]{chenevier}, since its determinant
over the residue field is that of the absolutely
irreducible representation over \(E\).
\proofstep{Recover matrices and flags from determinants.}

Thus a relative tangent gives a matrix cocycle modulo
infinitesimal conjugation. The polarization gives the transpose-inverse relation. For a fixed pairing matrix \(J\), an infinitesimal change \(1+\delta X\) with \(\delta^2=0\) preserves it exactly when
\[
 (1+\delta X)^{\mathsf T}J(1+\delta X)=J
 \quad\Longleftrightarrow\quad X^{\mathsf T}J+JX=0 .
\]
Normalization of the pairing uses division by two. The ordinary flags and their
character values are uniquely specified at first order
by the ordered eigenlines of a separating local
element. At each split \(\ell\)-place let \(\mu_v\) be the
preimage in the local Galois group of the \(\ell\)-primary
torsion in its topological abelianization. The ordered
characters \(\psi_{v,i}\) have prescribed values on
\(\mu_v\) from \(O^\times\). This is a condition in the
universal ordinary determinant ring, so an \(O\)-linear
tangent lift has the same values there, viewed as constants
in \(E[\delta]/(\delta^2)\).

For \(g\in\mu_v\) equality with the prescribed constant gives
\[
 \psi_{v,i}(g)(1+\delta c_{v,i}(g))=\psi_{v,i}(g)
 \quad\Longrightarrow\quad c_{v,i}(g)=0 .
\]
Here \(c_{v,i}(g)\in E\) and \(\psi_{v,i}(g)\) is a unit.
This condition survives reduction of an integral cocycle
modulo \(t^m\).
For the existence of the local flag, the trace relations and perfectness of the trace pairing give the ordered product identity
\[
 \prod_{i=1}^n\bigl(\rho(g_i)-\psi_i(g_i)I\bigr)=0\qquad(\text{all }g_1,\ldots,g_n).
\]
Take the same separating element at every
slot. Its distinct roots give the ordered spectral
projections, including over the dual numbers.
Replace only the \(k\)-th slot by \(g\). The factors
on the left kill the rows with index less than \(k\),
and those on the right kill the columns with index
greater than \(k\).

On the remaining rows and columns
they act by units. More explicitly, write \(L_k,R_k\) for the two products
of spectral factors at the separating element. For the
surviving rows and columns, the product equation reads
\[
 0=\bigl(L_k(\rho(g)-\psi_k(g))R_k\bigr)_{rc}
       =(L_k)_{rr}(\rho(g)-\psi_k(g))_{rc}(R_k)_{cc},
             \qquad r\ge k,\quad c\le k .
\]
Both outer factors are units by the distinct spectral values.
Taking \(k=c\) for each surviving column yields
\[
 \rho(g)_{rc}=0\ (r>c),\qquad \rho(g)_{cc}=\psi_c(g).
\]
Indeed the outer scalar factors are units, and the subtracted diagonal term occurs only when \(r=c\).
At the \(q\) new places simple Frobenius roots likewise
give the ordered lifts uniquely. The two tame and
Frobenius relations determine the new character
variables on those eigenlines. The determinant
coefficients and these parameters generate the chart.

Consequently the relative tangent injects into the
ordinary first cohomology of the limiting group with
its prescribed inertia conditions. For continuity of
the derivative, the kernel of evaluation onto \(K[[t]]\)
has a finite conormal module over \(K[[t]]\). The
derivative on it is bounded after multiplication by a
power of \(t\). It kills \(O[[s]]\), and if \(I\) denotes the evaluation kernel, the approximation by coefficient series reads
\[
 \mathfrak m_R^b\subseteq s^bO[[s]]+(I\cap\mathfrak m_R^b).
\]
Artin--Rees then gives a bound
tending to zero with \(b\). This applies to the
determinant coefficients and to the matrix coordinates
recovered by the inverse trace Gram matrix.

\proofstep{Local and global bound on relative tangents.}

For the dimension bound on this cohomology use the
first-order finite-level count. Over \(K[t]/(t^m)\)
it has the form
\[
 \bigl|\mathrm{Sel}_1(N,m)/\mathrm{Bd}_1(N,m)\bigr|
     \le |K|^{m(nd+nq)+c}.
\]
Here the bound holds for \(N\) sufficiently large for
each \(m\), with \(c\) independent of \(m,N\).
In the Euler-characteristic comparison let
\(\nu=n(n-1)/2\). At a split \(\ell\)-place, with
local degree \(d_v\) and principal-inertia rank
\(d'_v\), ordinary upper triangular cocycles contribute
at most
\[
 m\bigl(n(d'_v+1)+\nu(1+d_v)\bigr)+O(1)
\]
to the logarithmic count.

For the diagonal characters,
local class field theory identifies the free pro-\(\ell\)
part of inertia after quotienting by its torsion with
\(\mathbf Z_\ell^{d'_v}\). The vanishing on \(\mu_v\) leaves exactly \(d'_v\) inertia
parameters and one unramified parameter per character. Thus the
diagonals contribute logarithmic size at most \(mn(d'_v+1)\). For
upper entries filter the triangular cocycle space by
increasing distance from the diagonal. When all
preceding entries vanish, the \((i,j)\)-entry is a
cocycle with coefficients of character
\[
 \chi_{ij}=\psi_{v,i}/\psi_{v,j}.
\]

If \(\epsilon(g)/\chi_{ij}(g)-1\) has valuation \(b\)
at one local element, the invariants of the dual twist in
\(A/(t^m)\) lie in the kernel of multiplication by that scalar.
Over \(A/(t^m)\) the scalar generates the same ideal as \(t^b\). Consequently
\[
 \ker(t^b:A/(t^m)\to A/(t^m))=t^{\max(m-b,0)}A/(t^m),\qquad
 |\ker|=|K|^{\min(b,m)} .
\]
Writing \(D_{ij}\) for the rank-one coefficient module,
local duality, the Euler formula and the boundary size give
\[
 \begin{aligned}
 |Z^1(D_{ij})|
  &= |H^1(D_{ij})|\,|D_{ij}|/|H^0(D_{ij})|\\
  &= |D_{ij}|^{1+d_v}|H^0(D_{ij}^{\vee}(1))|
   \le |K|^{m(1+d_v)+b}.
 \end{aligned}
\]
Multiply these bounds at the successive strictly upper
positions of the filtration.

At the new places the unrestricted condition contributes
the invariants of the dual twist by local duality and
the tame Euler characteristic. For large \(N\) the
Frobenius matrix and cyclotomic value modulo \(t^m\)
are those of \(\sigma_j\). In the generic Frobenius eigenbasis the commutator has
multipliers \(a_i-a_j\) on off-diagonal entries. Hence
\[
     \dim_E\ker\operatorname{ad}(\rho(\sigma_j))=n.
\]
Only the diagonal entries contribute to the kernel rank. Smith normal form bounds the
size of the corresponding kernel modulo \(t^m\) by
\(|K|^{mn+O(1)}\).

The same argument bounds the term at
the old distinguished place by \(O(1)\) in logarithmic
size, using the vanishing local intertwining space.
The fixed-inertia primes give the unramified Euler
terms in the comparison. The global invariants
over \(E\) vanish. Schur's lemma on \(G_L\) makes an invariant \(X=\lambda I\). Invariance under complex conjugation then gives
\[
 \lambda I=-{}^t(\lambda I)=-\lambda I,\qquad 2\lambda=0 .
\]
The characteristic is odd, so \(X=0\). Finitely
many equations already have this kernel, so they
give a bounded invariant count at every truncation.

The dual auxiliary restrictions descend to the base
group, since they impose zero on the new tame generators,
and their Frobenius traces satisfy the chosen detection
tests. To bound this count modulo \(t^m\), work on the
topologically finitely generated group after the
pro-\(\ell\) quotient. The reduction of a cocycle also
factors through this quotient, by applying the
pro-\(\ell\) image argument to the cocycle and its
representation in their semidirect product.
A cocycle on the free profinite group with generators \(g_1,\ldots,g_m\) is determined by the tuple
\[
 z\longmapsto\bigl(z(g_1),\ldots,z(g_m)\bigr).
\]
Descent imposes
linear equations for relations. The rows of all these
equations and the finite trace tests span a finite
module of linear forms over \(A\), so finitely many
rows suffice even modulo \(t^m\).

Over \(E\) their
kernel consists of coboundaries by detection. Smith
normal form shows that the kernel modulo \(t^m\)
has size at most
\[
 |K|^{mr+O(1)}
\]
when its rank
over \(E\) is \(r\). For the coboundary map of the same generic rank \(r\), the
image has size at least \(|K|^{mr-O(1)}\). Thus division gives
\[
              |K|^{mr+O(1)}/|K|^{mr-O(1)}=|K|^{O(1)}.
\]
This bound is uniform in the truncation depth. These losses depend
on the fixed arc, the chosen local elements and the
finite trace tests. They are independent of \(N\) and
\(m\) in the indicated order of approximation.
Greenberg--Wiles duality subtracts
\[
 m\bigl((n+\nu)\#\{v\mid\ell\}+\nu d\bigr).
\]
With \(B_v\) the upper triangular matrix module at this
depth and \(A_v\) the full adjoint coefficient module, the
triangular coboundaries have order
\(|B_v|/|H^0(B_v)|\).

The ordinary local condition \(L_v\)
is an image of the triangular cocycle space, and dividing
by these coboundaries gives
\[
 \frac{|L_v|}{|H^0(A_v)|}
       \le \frac{|Z^1_{\rm ord}(B_v)|}{|B_v|}
                       \frac{|H^0(B_v)|}{|H^0(A_v)|}
       \le \frac{|Z^1_{\rm ord}(B_v)|}{|B_v|}.
\]
This is the subtraction of the triangular matrix size in
the Euler count.
At a real place the involution on the adjoint has
the form \(X\mapsto -S X^{\mathsf T} S^{-1}\) with \(S\)
symmetric and invertible. Since \(S^{\mathsf T}=S\), putting \(X=SY\) gives 
\[
 -SX^{\mathsf T}S^{-1}=X\quad\Longleftrightarrow\quad -Y^{\mathsf T}=Y.
\]
 This identifies the invariants with the skew matrices and accounts for \(\nu\) at each of the \(d\) real places.
Both sums of local degrees equal \(d\). The coefficient of
\(m\) after the global subtraction is
\[
 \sum_v\big(n(d'_v+1)+\nu(1+d_v)\big)
          +nq-(n+\nu)\#\{v\mid\ell\}-\nu d=nd+nq .
\]
This is the coefficient of the depth \(m\) in the first-order bound for one block. The local-condition comparison uses
the fixed torsion subgroup and the actual
decomposition maps at the finite levels.

For fixed \(m\), integral cocycles and the
relations at the limiting group are read at finite
levels on a set in \(\mathcal U\). Use finite generators and the
stable finite quotients in the common profinite group.
If \(a\) cocycles are independent over \(E\) modulo
boundaries, scale them by powers of \(t\) to make them
integral, and test their linear combinations against
coboundaries at finitely many elements. Such a finite
test suffices by Noetherianity of the module of linear
forms in the coefficients and the boundary matrix.

Smith normal form bounds the loss by a fixed power of
\(t\): if \(z_1,\ldots,z_a\) are these integral cocycles, then for \(m\ge c_0\)
\[
 \sum_j u_jz_j\equiv\text{a boundary}\pmod{t^m}\quad\Longrightarrow\quad(u_1,\ldots,u_a)\in t^{m-c_0}A^a.
\]
Thus a fibre of the map from the \(a\) coefficient tuples modulo \(t^m\) has size at most \(|K|^{a c_0}\), and the image has at least
\[
 |K|^{am}/|K|^{a c_0}=|K|^{a(m-c_0)}
\]
elements. The boundary matrices,
including their skew condition, lift from the finite
quotients. The cardinality estimate gives
\[
        a(m-c_0)\le m(nd+nq)+c \qquad(m\ge c_0).
\]
Dividing by \(m\) and taking the limit proves \(a\le nd+nq\). In this one-block case
the odd matrix module of the two-block comparison
is zero. Its second-order obstruction equations have
no odd entries. Only the first-order dimension bound
enters the cotangent argument.

By completeness and the relative cotangent bound,
\(C/IC\) is a quotient of a power-series ring in
\(D+nq\) variables over the coefficient DVR. Lift a
spanning set of the relative cotangent and use
successive approximation to obtain the surjection.
The power-series ring is a domain of dimension
\(1+D+nq\). Let \(S\) denote this power series ring and \(P\) the inverse image of \(Z/IC\). Its full-dimensional quotient forces
\[
 \operatorname{ht}P+\dim(S/P)\le\dim S=\dim(S/P),
 \qquad \operatorname{ht}P=0.
\]
Since \(S\) is a domain, \(P=0\), and the kernel of \(S\to C/IC\) lies in \(P\). Hence \(C/IC\) itself is a domain and
\(Z=IC\). Every prime containing \(IC\) now lies in
the support of \(P_C\).

\proofstep{Descend the vanishing to the original Hecke quotient.}

The patched calculation supplies support after the auxiliary choices and coefficient extensions. The point on the original Hecke quotient is the one needed in the propagation argument.


Finally, a prime below \(p\) lifts below the chosen
arc prime of \(R_0'\) by flat going down. Its inverse
image in \(R\) contains \(I\) and \(\mathfrak y\).
Lift this prime to \(C\) by faithful flatness of
completion of \(R_{p_R}\). The augmentation support
comparison puts \(J_0\) in its contraction to \(R_0\).
This proves the one-point implication from the
finite-level module and cohomological inputs above.
\label{cor:ordinary-ordinary-patching-1}
\end{proof}
